\section*{2.5 Machine Maintenance}

This is a finite-horizon stochastic DP problem. The state is the condition of the machine at the start of a week:
running \(R\) or broken \(B\). The action is chosen after observing this state,
so the solution is a closed-loop policy rather than one fixed open-loop action
sequence.

Let \(V_R(h)\) and \(V_B(h)\) be the optimal expected profit with \(h\) weeks
remaining when the machine starts the week running or broken. Since this
problem asks for profit, the Bellman recursion uses a maximum:
\[
V_h(x)=\max_{u\in U(x)}\mathbb E\left[
    \text{one-week profit}+V_{h-1}(x')
\right]
\]

For a running machine, the two choices are preventive maintenance \(M\) and no
maintenance \(N\). The one-step action values are
\[
\begin{aligned}
Q_R^M(h)&=0.6\left(80+V_R(h-1)\right)
        +0.4\left(-20+V_B(h-1)\right)\\
Q_R^N(h)&=0.3\left(100+V_R(h-1)\right)
        +0.7V_B(h-1)
\end{aligned}
\]
Here \(80=100-20\) is the profit if maintenance is done and the machine runs,
and \(-20\) is the maintenance cost if it still breaks.

For a broken machine, the two choices are repair \(C\) and replacement \(P\).
The corresponding action values are
\[
\begin{aligned}
Q_B^C(h)&=0.6\left(60+V_R(h-1)\right)
        +0.4\left(-40+V_B(h-1)\right)\\
Q_B^P(h)&=-50+V_R(h-1)
\end{aligned}
\]
The repair value uses \(60=100-40\), while replacement gives a new machine
that is guaranteed to run for the current week, so its immediate net profit is
\(100-150=-50\).

The DP update is therefore
\[
\begin{aligned}
V_R(h)&=\max\left\{
40+0.6V_R(h-1)+0.4V_B(h-1),\
30+0.3V_R(h-1)+0.7V_B(h-1)
\right\}\\
V_B(h)&=\max\left\{
20+0.6V_R(h-1)+0.4V_B(h-1),\
-50+V_R(h-1)
\right\}
\end{aligned}
\]
with terminal values \(V_R(0)=V_B(0)=0\).

Working backward gives the action values
\[
\begin{array}{c|cc|cc}
h & Q_R^M(h) & Q_R^N(h) & Q_B^C(h) & Q_B^P(h)\\
\hline
1 & 40 & 30 & 20 & -50\\
2 & 72 & 56 & 52 & -10\\
3 & 104 & 88 & 84 & 22
\end{array}
\]
so the value functions are
\[
\begin{array}{c|cccc}
h&0&1&2&3\\
\hline
V_R(h)&0&40&72&104\\
V_B(h)&0&20&52&84
\end{array}
\]
At each remaining horizon, preventive maintenance is better when the machine is
running, and repair is better when the machine is broken.

The first week is handled separately because the machine is new and guaranteed
to run through that week. After receiving the first week's profit, there are
three weeks left with the machine running, so the total expected profit is
\[
    \boxed{100+V_R(3)=204}
\]
